\documentclass[a4paper,oneside,11pt]{book}
\usepackage{graphicx}
\usepackage{titling}
\usepackage{geometry}
\usepackage{listings}
\usepackage{xcolor}
\usepackage{float}
\usepackage{indentfirst}
\usepackage{longtable}
\usepackage[T1]{fontenc}
\usepackage{lmodern}
\usepackage[hyphens]{url}
\usepackage{hyperref}
\def\UrlBreaks{\do\/\do-}

\geometry{left=3cm, right=2.5cm, top=3cm, bottom=3cm}
\emergencystretch=5em

\definecolor{codegreen}{rgb}{0,0.6,0}
\definecolor{codegray}{rgb}{0.5,0.5,0.5}
\definecolor{codepurple}{rgb}{0.58,0,0.82}
\definecolor{backcolour}{rgb}{0.97,0.97,0.95}
\definecolor{codeblue}{rgb}{0.16,0.32,0.75}

\lstdefinelanguage{yaml}{
    keywords={true,false,null,yes,no,on,needs,if,runs-on,uses,with,run,steps,jobs,name,env,environment,cache},
    keywordstyle=\color{codeblue}\bfseries,
    basicstyle=\ttfamily\footnotesize,
    sensitive=false,
    comment=[l]{\#},
    commentstyle=\color{codegray}\itshape,
    stringstyle=\color{codepurple},
    morestring=[b]',
    morestring=[b]"
}

\lstdefinestyle{codestyle}{
    backgroundcolor=\color{backcolour},
    basicstyle=\ttfamily\footnotesize,
    commentstyle=\color{codegreen},
    keywordstyle=\color{codeblue}\bfseries,
    stringstyle=\color{codepurple},
    numberstyle=\tiny\color{codegray},
    numbers=left,
    numbersep=6pt,
    showspaces=false,
    showstringspaces=false,
    breaklines=true,
    tabsize=2,
    frame=single,
    rulecolor=\color{gray},
    captionpos=b,
}
\lstset{style=codestyle}

\title{Laporan Praktikum \\ Konstruksi dan Evolusi Perangkat Lunak\\
Pertemuan 4 \\ Frontend Vue 3 dan Pipeline CI/CD Fullstack}
\author{Hisyam Al Rayyan / 24/541582/SV/24926}
\date{\today}

\begin{document}

\begin{titlingpage}
\begin{center}
\vspace{4cm}
\begin{huge}
\textbf{\thetitle} \\
\end{huge}
\vspace{2cm}
\includegraphics[height=8cm]{images/logo_ugm.png}\\
\vspace{4cm}
{\Large
\begin{tabular}{rl}
Nama            & : Hisyam Al Rayyan \\
NIM             & : 24/541582/SV/24926 \\
Kelas           & : A2 \\
Dosen Pengampu  & : Galih Malela Damaraji \\
\end{tabular}
}
\vspace{2cm}
\end{center}
\end{titlingpage}

\tableofcontents

%=== BAB 1 ===
\chapter{Pendahuluan}

\section{Tujuan Praktikum}
\begin{enumerate}
    \item Membuat endpoint JSON dari tabel CRUD Laravel yang dapat diakses oleh aplikasi frontend.
    \item Membangun aplikasi Vue 3 dengan Vue Router minimal dua halaman.
    \item Mengambil URL API dari variabel lingkungan \texttt{VITE\_API\_URL} tanpa menulis URL langsung di kode.
    \item Menyusun workflow CI/CD empat job untuk frontend menggunakan \texttt{npm ci} dan cache npm.
    \item Meneruskan artefak hasil \texttt{npm run build} ke job deploy tanpa menjalankan build ulang.
    \item Menulis minimal satu unit test Vitest yang berjalan tanpa Laravel.
    \item Membuktikan pipeline berhenti ketika test gagal dan deploy dilewati pada Pull Request.
\end{enumerate}

\section{Alat dan Bahan}
\begin{enumerate}
    \item Laptop dengan sistem operasi Windows.
    \item PHP 8.4, Composer, Node.js, dan npm.
    \item Framework Laravel 13 sebagai backend.
    \item Vue 3 dengan Vite sebagai frontend.
    \item Vitest sebagai framework pengujian unit frontend.
    \item GitHub Actions sebagai platform CI/CD.
\end{enumerate}

%=== BAB 2 ===
\chapter{Hasil dan Pembahasan}

\section{Endpoint JSON Laravel}

Endpoint JSON disediakan dari tabel \texttt{tugas} yang telah dibuat pada praktikum sebelumnya. Endpoint \texttt{GET /api/tugas} mengembalikan data dalam format JSON. Untuk mendukung akses dari aplikasi Vue yang berjalan di port berbeda, middleware CORS dikonfigurasi pada Laravel.

\begin{lstlisting}[language=php, caption={Endpoint API pada routes/api.php}]
<?php

use App\Http\Controllers\TugasController;
use Illuminate\Support\Facades\Route;

Route::get('/tugas', [TugasController::class, 'index']);
Route::post('/tugas', [TugasController::class, 'store']);
Route::get('/tugas/{id}', [TugasController::class, 'show']);
\end{lstlisting}

% Screenshot 1: Endpoint JSON berjalan
%\begin{figure}[H]
%    \centering
%    \includegraphics[width=0.8\textwidth]{4-1.png}
%    \caption{Endpoint GET /api/tugas mengembalikan data JSON}
%    \label{fig:4-1}
%\end{figure}

\section{Aplikasi Vue 3}

Aplikasi Vue 3 dibuat di folder \texttt{frontend/} menggunakan Vite. Aplikasi memiliki dua halaman yang diatur dengan Vue Router:
\begin{itemize}
    \item Halaman \textbf{Home} (\texttt{/}): Menampilkan daftar data tugas yang diambil dari Laravel API.
    \item Halaman \textbf{About} (\texttt{/about}): Menampilkan informasi aplikasi.
\end{itemize}

URL API diambil dari variabel lingkungan \texttt{VITE\_API\_URL} menggunakan \texttt{import.meta.env.VITE\_API\_URL}, tidak ditulis langsung di kode.

\begin{lstlisting}[language=javascript, caption={Pengambilan data dari API menggunakan VITE\_API\_URL}]
// src/views/HomeView.vue
const apiUrl = import.meta.env.VITE_API_URL
const response = await fetch(`${apiUrl}/api/tugas`)
const data = await response.json()
\end{lstlisting}

% Screenshot 2: Halaman Home Vue
%\begin{figure}[H]
%    \centering
%    \includegraphics[width=0.8\textwidth]{4-2.png}
%    \caption{Halaman Home menampilkan data dari Laravel API}
%    \label{fig:4-2}
%\end{figure}

% Screenshot 3: Halaman About Vue
%\begin{figure}[H]
%    \centering
%    \includegraphics[width=0.8\textwidth]{4-3.png}
%    \caption{Halaman About aplikasi Vue}
%    \label{fig:4-3}
%\end{figure}

\section{Unit Test Vitest}

Minimal satu unit test ditulis menggunakan Vitest untuk menguji logika aplikasi Vue tanpa memerlukan Laravel berjalan. Test memverifikasi fungsi pemformatan data yang digunakan di halaman Home.

\begin{lstlisting}[language=javascript, caption={Unit test Vitest untuk logika frontend}]
// src/__tests__/tugas.test.js
import { describe, it, expect } from 'vitest'
import { formatStatus } from '../utils/tugas'

describe('formatStatus', () => {
  it('returns label Selesai for status selesai', () => {
    expect(formatStatus('selesai')).toBe('Selesai')
  })

  it('returns label Pending for status pending', () => {
    expect(formatStatus('pending')).toBe('Pending')
  })
})
\end{lstlisting}

% Screenshot 4: Hasil Vitest
%\begin{figure}[H]
%    \centering
%    \includegraphics[width=0.8\textwidth]{4-4.png}
%    \caption{Hasil eksekusi unit test Vitest di GitHub Actions}
%    \label{fig:4-4}
%\end{figure}

\section{Workflow CI/CD Frontend}

Workflow CI/CD untuk frontend dikonfigurasi pada \texttt{.github/workflows/frontend.yml} dengan empat job yang dirangkai menggunakan \texttt{needs:}.

\begin{lstlisting}[language=yaml, caption={Workflow frontend empat job}]
name: Frontend CI/CD

on:
  push:
    branches: ['**']
  pull_request:
    branches: [main, dev]

jobs:
  lint:
    runs-on: ubuntu-latest
    defaults:
      run:
        working-directory: frontend
    steps:
      - uses: actions/checkout@v4
      - uses: actions/setup-node@v4
        with:
          node-version: '20'
          cache: 'npm'
          cache-dependency-path: frontend/package-lock.json
      - run: npm ci
      - run: npm run lint

  test:
    runs-on: ubuntu-latest
    needs: lint
    defaults:
      run:
        working-directory: frontend
    steps:
      - uses: actions/checkout@v4
      - uses: actions/setup-node@v4
        with:
          node-version: '20'
          cache: 'npm'
          cache-dependency-path: frontend/package-lock.json
      - run: npm ci
      - run: npm run test

  build:
    runs-on: ubuntu-latest
    needs: test
    defaults:
      run:
        working-directory: frontend
    steps:
      - uses: actions/checkout@v4
      - uses: actions/setup-node@v4
        with:
          node-version: '20'
          cache: 'npm'
          cache-dependency-path: frontend/package-lock.json
      - run: npm ci
      - run: npm run build
        env:
          VITE_API_URL: ${{ vars.API_URL }}
      - uses: actions/upload-artifact@v4
        with:
          name: frontend-dist
          path: frontend/dist/

  deploy:
    runs-on: ubuntu-latest
    needs: build
    if: github.ref == 'refs/heads/main'
    environment: production
    steps:
      - name: Download build artifact
        uses: actions/download-artifact@v4
        with:
          name: frontend-dist
          path: dist/
      - name: Show dist contents
        run: ls -la dist/
      - name: Deploy frontend
        run: |
          echo "=== FRONTEND DEPLOY ==="
          echo "Deploying files from dist/ to server..."
          echo "Deploy complete!"
\end{lstlisting}

% Screenshot 5: Pipeline frontend
%\begin{figure}[H]
%    \centering
%    \includegraphics[width=0.8\textwidth]{4-5.png}
%    \caption{Pipeline CI/CD frontend empat job di GitHub Actions}
%    \label{fig:4-5}
%\end{figure}

\section{Pembuktian: Pull Request Deploy Skipped}

Untuk membuktikan bahwa job \texttt{deploy} dilewati pada Pull Request, dibuat Pull Request dari branch fitur. Pipeline menjalankan \texttt{lint}, \texttt{test}, dan \texttt{build}, sedangkan job \texttt{deploy} berstatus \textit{skipped}.

% Screenshot 6: PR pipeline deploy skipped
%\begin{figure}[H]
%    \centering
%    \includegraphics[width=0.8\textwidth]{4-6.png}
%    \caption{Pipeline pada Pull Request: deploy skipped}
%    \label{fig:4-6}
%\end{figure}

\section{Pembuktian: Test Gagal Menghentikan Pipeline}

Satu assertion pada unit test Vitest sengaja diubah menjadi salah. Setelah push, pipeline menunjukkan job \texttt{test} berstatus merah dan job selanjutnya tidak dijalankan. Setelah perbaikan, pipeline kembali hijau seluruhnya.

% Screenshot 7: Pipeline merah test gagal
%\begin{figure}[H]
%    \centering
%    \includegraphics[width=0.8\textwidth]{4-7.png}
%    \caption{Pipeline frontend berhenti merah karena unit test gagal}
%    \label{fig:4-7}
%\end{figure}

% Screenshot 8: Pipeline hijau setelah perbaikan
%\begin{figure}[H]
%    \centering
%    \includegraphics[width=0.8\textwidth]{4-8.png}
%    \caption{Pipeline frontend hijau setelah test diperbaiki}
%    \label{fig:4-8}
%\end{figure}

%=== BAB 3 ===
\chapter{Kesimpulan}

Berdasarkan praktikum yang telah dilakukan, dapat disimpulkan beberapa hal berikut:
\begin{enumerate}
    \item Endpoint JSON \texttt{GET /api/tugas} berhasil disediakan oleh Laravel dan dapat diakses oleh aplikasi Vue 3 menggunakan konfigurasi CORS yang tepat.
    \item Aplikasi Vue 3 berhasil dibuat dengan dua halaman menggunakan Vue Router. URL API diambil dari \texttt{VITE\_API\_URL} tanpa penulisan langsung di kode.
    \item Workflow CI/CD empat job (\texttt{lint}, \texttt{test}, \texttt{build}, \texttt{deploy}) berhasil dikonfigurasi dengan \texttt{npm ci} dan cache npm untuk efisiensi.
    \item Artefak hasil \texttt{npm run build} berhasil diteruskan ke job \texttt{deploy} tanpa menjalankan build ulang, menggunakan mekanisme \texttt{upload-artifact} dan \texttt{download-artifact}.
    \item Unit test Vitest berhasil berjalan di GitHub Actions tanpa memerlukan Laravel berjalan.
    \item Pipeline terbukti berhenti ketika unit test gagal, dan job \texttt{deploy} terbukti dilewati pada Pull Request.
\end{enumerate}

\clearpage

\begin{thebibliography}{9}
\bibitem{vue3Doc}
Vue.js. (2026). \textit{Vue 3 Documentation}. Diakses dari \url{https://vuejs.org/guide}.
\bibitem{vitestDoc}
Vitest. (2026). \textit{Vitest Documentation}. Diakses dari \url{https://vitest.dev/guide}.
\bibitem{githubActionsDoc}
GitHub. (2026). \textit{GitHub Actions Documentation}. Diakses dari \url{https://docs.github.com/en/actions}.
\end{thebibliography}

\end{document}
